\subsection{Ordinary eighth moments for odd terminals with an inward arm of length at least two}\label{proof:Q5-C093}

\paragraph{Q93.0 Structures and directions.}
After maximal equal pairing, let the distinct residual arms form
\begin{equation}\tag{Q93.0.1}\label{eq:Q93:0.1}
a<b_1<c_1<\cdots<b_k<c_k,\qquad a\ge2,\quad k\ge1.
\end{equation}
Every edge of the \(a\)-arm points toward the terminal. Each \(b_i\)-arm
has only its first edge inward and all remaining edges outward;
each \(c_i\)-arm points outward. All connector edges point from the
terminal to the receiving root, which has an outward unit arm. With
genuine equal-pair products \(E_t,E_r\), set
\begin{equation}\tag{Q93.0.2}\label{eq:Q93:0.2}
F=\bar f_a\prod_{i=1}^kh_{b_i}f_{c_i},\qquad
X_D=T^L(E_tF),\qquad H_D=E_rpX_D.
\end{equation}
These directions depend only on the arm list. For every \(L\ge1\),
all finite lengths and gaps, and all genuine decorations, we prove
\begin{equation}\tag{Q93.0.3}\label{eq:Q93:0.3}
\boxed{\E H_D^8\le Q(d)\le e^{-d}\le1-3d/4,}
\end{equation}
The three ranges and bounds are listed in Table~\ref{tab:oc:q93-ranges}.
\begin{table}[htbp]
\centering\small
\setlength{\tabcolsep}{5pt}
\caption{The three ordinary eighth-moment ranges in \eqref{eq:Q93:0.3}.}\label{tab:oc:q93-ranges}
\begin{tabular}{lcl}
\toprule
First short arm & Pair count & $Q(d)$\\
\midrule
$b_1=3,\ a=2$ & $k\ge3$ & $e^{-[10+(128/5)(k-3)]d}$\\
$b_1=4,\ a\in\{2,3\}$ & $k\ge2$ & $e^{-[17+(128/5)(k-2)]d}$\\
$b_1\ge5,\ 2\le a<b_1$ & $k\ge1$ & $e^{-[1+(128/5)(k-1)]d}$\\
\bottomrule
\end{tabular}
\end{table}


\paragraph{Q93.1 The finite calculus.}
Use the fixed finite probability space and operators of
Appendix~\ref{proof:Q5-C064}, with
\begin{equation}\tag{Q93.1.1}\label{eq:Q93:1.1}
\begin{gathered}
Jv=(\E v)\one,\quad A=J+K,\quad T=J-K=A^*,\quad
g=K\one,\quad p=1+g,\quad q=1-g,\\
f_n=A^n\one,\quad\bar f_n=T^n\one,\quad h_n=Tf_{n-1},\quad
\delta_i=\|K^i\one\|_2^2,\quad d=\delta_1,\\
\theta=13/32,\quad\kappa=19/32,\quad\tau=d-\delta_2.
\end{gathered}
\end{equation}
The finite skew and path identities give
\begin{equation}\tag{Q93.1.2}\label{eq:Q93:1.2}
\begin{gathered}
0\le d\le1/3,\quad d^2\le\delta_2\le\theta d,\quad
\delta_{i+1}\le\theta\delta_i,\quad\|K\|^2\le\theta,\quad
|Kv(x)|\le\|v\|_2,\\
0\le f_n,\bar f_n,h_n\le2,\qquad
\mu_n:=\E f_n=\E\bar f_n\in[0,1],\\
\E K^{2i}\one=(-1)^i\delta_i,\qquad\E K^{2i+1}\one=0.
\end{gathered}
\end{equation}
After zero-mass atoms are removed, \(p(x),q(x)\ge\mu(x)>0\).
Each tree message is positive: conditional on the root label, the
all-equal labelling has positive mass and all diagonal edge factors
are one. At \(d=0\), \(g=0\) and leaf integration gives every message
equal to one. We henceforth take \(d>0\) where division is used.
For every genuine equal pair,
\begin{equation}\tag{Q93.1.3}\label{eq:Q93:1.3}
f_\ell+h_\ell=2\mu_{\ell-1},\qquad
0\le f_\ell h_\ell\le\mu_{\ell-1}^2\le1.
\end{equation}
With \(\beta_1=\beta_2=0\),
\begin{equation}\tag{Q93.1.4}\label{eq:Q93:1.4}
\begin{gathered}
\beta_n=\sum_{i=1}^{\lfloor(n-1)/2\rfloor}(-1)^{i-1}\delta_i
\quad(n\ge3),\qquad G_n(z)=\sum_{i=0}^{n-1}z^i,\\
S_n=\sum_{i=0}^nK^i\one=f_n+\sum_{i=3}^n\beta_i f_{n-i},
\quad\kappa d\le\tau\le\beta_n\le d\quad(n\ge3),\\
\beta_3=\beta_4=d,\quad
\|G_n(K)\|\le1\quad(n=0,1\text{ or }n\ge3),\quad
\|G_2(K)\|\le\sqrt{1+\theta}<6/5.
\end{gathered}
\end{equation}
Under reversal put \(K'=-K\), \(g'=-g\), and
\(S'_n=\sum_{i=0}^nK'^i\one\); the same convolution uses \(\bar f_n\).
Appendix~\ref{proof:Q5-C081} proves, for every real \(V\),
\begin{equation}\tag{Q93.1.5}\label{eq:Q93:1.5}
\boxed{\E e^{8K'V}\le e^{8\E(K'V)+17\|V\|_2^2}.}
\end{equation}

\paragraph{Q93.2.1 Pair losses at every arm length.}
Let \(R_{st}=\sum_{r=s+1}^tK^r\one\), \(s<t\).
The finite difference formula gives, for \(s=3\),
\begin{equation}\tag{Q93.2.1}\label{eq:Q93:2.1}
h_3+f_t\le2+R_{3t}-(2d-\delta_2)
\le2+R_{3t}-(2-\theta)d.
\end{equation}
For \(s=4\), retain the negative convolution term:
\fitdisplay{
f_t-f_4\le R_{4t}+\delta_2(1+p)-\tau f_2,\qquad f_2\ge p^2/2.
}{}
Completing the square in \(p\ge0\), with \(\delta_2\le d/2\) and
\(\tau\ge d/2\), gives
\(\delta_2(1+p)-\tau f_2\le\delta_2+\delta_2^2/(2\tau)\le3d/4\).
Since \(h_4+f_4=2-4d\), this yields
\(h_4+f_t\le2+R_{4t}-13d/4\), including \(t=5\).
For \(s\ge5\), the all-index mean proof in
Appendix~\ref{proof:Q5-C090} gives \(\mu_n\le1-83d/32\) for \(n\ge4\).
It treats \(n=4,5,6,7\) explicitly and retains \(f_0,\ldots,f_5\) and
three higher losses in the convolution for \(n\ge8\), giving loss at
least \(95d/36\). The positive part of the finite difference is at
most \(4\delta_2/(1-\theta^2)\le1664d/855\), so the pair-sum loss is
\fitdisplay{
83d/16-1664d/855=(16/5+113/2736)d>16d/5.
}{}
The sums are nonnegative. AM--GM and \(1+u\le e^u\) give
\begin{equation}\tag{Q93.2.2}\label{eq:Q93:2.2}
h_sf_t\le
\begin{cases}
e^{R_{3t}-(51/32)d},&s=3,\\
e^{R_{4t}-(13/4)d},&s=4,\\
e^{R_{st}-(16/5)d},&s\ge5.
\end{cases}
\end{equation}
Later pairs of minimum at least four may use the weaker loss \(16d/5\).
Thus
\begin{equation}\tag{Q93.2.3}\label{eq:Q93:2.3}
\begin{gathered}
F\le\bar f_ae^{R-\Gamma},\qquad
R=\sum_{i=1}^k\sum_{r=b_i+1}^{c_i}K^r\one,\\
\Gamma/d=\begin{cases}
51/32+(16/5)(k-1),&b_1=3,\\
13/4+(16/5)(k-1),&b_1=4,\\
16k/5,&b_1\ge5.
\end{cases}
\end{gathered}
\end{equation}

\paragraph{Q93.2.2 Finite masks and the exponential remainder.}
The intervals are disjoint, so \(R=K^{b_1}P(K)g\), where \(P\) is a
finite zero-one polynomial with arbitrary gaps. The finite mask proof
in \eqref{eq:L90:masknorm} gives, at \(z=i\sqrt x\), \(0\le x\le\theta\),
\begin{equation}\tag{Q93.2.4}\label{eq:Q93:2.4}
|P(z)|^2\le\frac{1+x}{(1-x^2)^2}
\le C_*:=\frac{32768}{16245}<\frac{100}{49}.
\end{equation}
The spectral and row bounds respectively imply
\begin{equation}\tag{Q93.2.5}\label{eq:Q93:2.5}
\|R\|^2\le\theta^{b_1}C_*d,\qquad
|R|^2\le\theta^{b_1-1}C_*d.
\end{equation}
Since \(b_1\ge3\) and \(\theta^2C_*=5408/16245<1/3\),
\(|R|^2<d/3\le1/9\). The factorial bound
\(n!\ge2\cdot3^{n-2}\) bounds the exponential remainder by \(9R^2/16\)
for \(|R|\le1/3\). Hence throughout the domain
\begin{equation}\tag{Q93.2.6}\label{eq:Q93:2.6}
e^R\le1+R+3R^2/5.
\end{equation}
Use the scalar norm constants
\begin{equation}\tag{Q93.2.7}\label{eq:Q93:2.7}
\|R\|\le\varrho\sqrt d,\qquad
\varrho=3/8\ (b_1=3),\quad1/4\ (b_1=4),\quad3/20\ (b_1\ge5).
\end{equation}
Indeed,
\[
\begin{gathered}
\theta^3C_*=2197/16245<9/64,\\
\theta^4C_*=28561/519840<1/16,\\
\theta^5C_*=371293/16634880<9/400.
\end{gathered}
\]

\paragraph{Q93.2.3 Mean and centered part of the inward arm.}
Put \(\mu=\E\bar f_a=\mu_a\in[0,1]\), \(v=\bar f_a-\mu\),
\(\Pi=I-J\), and \(C_0=\Pi K\Pi\). The centered path identity is
\begin{equation}\tag{Q93.2.8}\label{eq:Q93:2.8}
\bar f_a-\mu_a=-\sum_{i=0}^{a-1}\mu_{a-1-i}(-C_0)^ig.
\end{equation}
Every coefficient is in \([0,1]\). Applying the same finite mask
spectral bound to the skew compression gives
\(\sigma(\bar f_a)\le\sqrt{C_*}\sqrt d<(3/2)\sqrt d\).
For \(b_1=3,4\), only \(a=2,3\) are possible. Direct expansion gives
\(\sigma(\bar f_2)^2=d+\delta_2-d^2\le36d/25\).
For \(\bar f_3\), its centered expression
\(-(g+K^3\one)+(Kg+d)\) has orthogonal summands, so
\(\sigma(\bar f_3)^2=d-\delta_2+\delta_3-d^2\le d\).
Consequently
\begin{equation}\tag{Q93.2.9}\label{eq:Q93:2.9}
\sigma(\bar f_a)\le s\sqrt d,\qquad
s=6/5\ (b_1=3,4),\quad s=3/2\ (b_1\ge5).
\end{equation}

\paragraph{Q93.3 Positive parts of all finite propagation means.}
Let \(D_i=\E K'^iR\) and write
\(R=\sum_{r\ge b_1+1}c_rK^r\one\), a finite sum with \(c_r\in\{0,1\}\).
A term with odd \(i+r\) has zero mean; otherwise its contribution is
\((-1)^{i+(i+r)/2}c_r\delta_{(i+r)/2}\).
Positive terms can therefore occur precisely when \(r\equiv i\pmod4\).
For \(i=1,2,3,4\), their first possible moments are
\begin{table}[htbp]
\centering\small
\setlength{\tabcolsep}{5pt}
\caption{First possible positive moments in $D_i=\E B^iR$. Adding four to either index raises the moment index by two.}\label{tab:oc:q93-positive-starts}
\begin{tabular}{ccccc}
\toprule
$b_1$ & $i=1$ & $i=2$ & $i=3$ & $i=4$\\
\midrule
$3$ & $\delta_3$ & $\delta_4$ & $\delta_5$ & $\delta_4$\\
$4$ & $\delta_3$ & $\delta_4$ & $\delta_5$ & $\delta_6$\\
$\ge5$ & $\delta_5$ & $\delta_4$ & $\delta_5$ & $\delta_6$\\
\bottomrule
\end{tabular}
\end{table}
Adding four to either index increases the moment index by two.
Bounding both finite tails by their geometric upper sums yields
\begin{equation}\tag{Q93.3.1}\label{eq:Q93:3.1}
\sum_{i\ge1}(D_i)_+\le
\begin{cases}
\dfrac{\theta^2}{(1-\theta)^2}d=\dfrac{169}{361}d<d/2,&b_1=3,\\[1mm]
\dfrac{\theta^2(1+\theta+\theta^2+\theta^3)}{(1-\theta^2)^2}d
=\dfrac{201617}{519840}d<2d/5,&b_1=4,\\[1mm]
\dfrac{\theta^3}{(1-\theta)^2}d=\dfrac{2197}{11552}d<d/5,&b_1\ge5.
\end{cases}
\end{equation}
The bound applies to each finite initial sum used by Duhamel.
The same sign calculation gives
\begin{table}[htbp]
\centering\small
\setlength{\tabcolsep}{5pt}
\caption{Upper bounds for individual propagation means.}\label{tab:oc:q93-individual-means}
\begin{tabular}{cccc}
\toprule
$b_1$ & $D_0$ & $D_1$ & $D_i,\ i\ge2$\\
\midrule
$3$ & $d/2$ & $d/5$ & $d/12$\\
$4$ & $d/12$ & $d/5$ & $d/12$\\
$\ge5$ & $d/12$ & $d/30$ & $d/12$\\
\bottomrule
\end{tabular}
\end{table}
For \(D_0\), positive terms start at \(\delta_2\) or \(\delta_4\);
for \(D_1\), at \(\delta_3\) or \(\delta_5\); for \(i\ge2\), their
first possible moment index is at least four. The needed bounds are
\fitdisplay{
\begin{gathered}
\theta/(1-\theta^2)=416/855<1/2,\quad
\theta^2/(1-\theta^2)=169/855<1/5,\\
\theta^3/(1-\theta^2)=2197/27360<1/12,\quad
\theta^4/(1-\theta^2)=28561/875520<1/30.
\end{gathered}
}{}
Deleting an absent positive term only improves these upper bounds.

\paragraph{Q93.4 The common inward-arm propagation lemma.}
Apply Lemma~\ref{lem:oc:anchored-eighth} with $B=K'=-K$,
$\mu=\E\bar f_a$, and the values of $s,\varrho$ in
\eqref{eq:Q93:2.7} and \eqref{eq:Q93:2.9}. Its first-step constant
$c_0=2s\varrho+(12/5)\varrho^2$ equals $99/80$, $3/4$, and
$63/125$ in the three rows. The individual mean bounds in
Table~\ref{tab:oc:q93-individual-means} therefore give
\begin{equation}\tag{Q93.4.2}\label{eq:Q93:4.2}
 \delta_+\le
 \begin{cases}7d/4,&b_1=3,\\5d/6,&b_1=4,\\3d/5,&b_1\ge5.
 \end{cases}
\end{equation}
For the last row, $(D_0)_++63d/125\le881d/1500<3d/5$.
The lemma retains the coefficient $\mu$ of every perturbation
and gives the exact expressions \eqref{eq:Q93:4.3}--\eqref{eq:Q93:4.6}.

\paragraph{Q93.5 Every nonunit connector.}
For \(L\ge2\), \(n\ge4\), so \(G_{n-1}\) is a contraction.
Using \(\mu\le1\) and \(\sqrt\theta<13/20\),
\begin{equation}\tag{Q93.5.1}\label{eq:Q93:5.1}
\|V\|\le(1+\sqrt\theta\,\varrho)\sqrt d.
\end{equation}
The three coefficients are smaller than \(5/4,7/6,11/10\);
the middle one is \(1+13/80=93/80<7/6\).
Equations \eqref{eq:Q93:3.1}, \eqref{eq:Q93:4.2} and
\eqref{eq:Q93:4.5} give the complete cost table
\begin{table}[htbp]
\centering\small
\setlength{\tabcolsep}{5pt}
\caption{Nonunit connector costs in \eqref{eq:Q93:4.6}. The last column excludes $-8\Gamma$. All entries are upper bounds.}\label{tab:oc:q93-nonunit-cost}
\begin{tabular}{ccccc}
\toprule
$b_1$ & $c_L/d$ & $\E N/d$ & $\|V\|^2/d$ & Cost divided by $d$\\
\midrule
$3$ & $9/2-\kappa$ & $-\kappa+1/12$ & $25/16$ & $2579/48$\\
$4$ & $37/15-\kappa$ & $-\kappa+1/12$ & $49/36$ & $6127/180$\\
$\ge5$ & $8/5-\kappa$ & $-\kappa+1/12$ & $121/100$ & $119/30+2057/100$\\
\bottomrule
\end{tabular}
\end{table}
The first row consists of costs \(163d/6\) and \(425d/16\), totaling
\(2579d/48\); the second, \(109d/10\) and \(833d/36\), totaling
\(6127d/180\); the third, \(119d/30\) and \(2057d/100\).
At the minimum allowed pair counts, the net losses are
\begin{equation}\tag{Q93.5.2}\label{eq:Q93:5.2}
\begin{aligned}
b_1=3, k=3:&\quad1279d/20-2579d/48=2453d/240>10d,\\
b_1=4, k=2:&\quad258d/5-6127d/180=3161d/180>17d,\\
b_1\ge5, k=1:&\quad128d/5-119d/30-2057d/100=319d/300>d.
\end{aligned}
\end{equation}
Each additional pair adds \(16d/5\) to \(\Gamma\) without increasing
the mask costs, and therefore adds \(128d/5\) to the eighth-moment loss.

\paragraph{Q93.6 Unit connectors and the \(G_2\) exception.}
For \(L=1\), retain the exact value \(\bar f_0=1\). Then
\begin{equation}\tag{Q93.6.1}\label{eq:Q93:6.1}
c_1+\E N\le\delta+\mu D_1-\beta_{a+1}-\beta_{a+2},\qquad
\|V\|\le(\|G_a\|+\varrho)\sqrt d.
\end{equation}
For $a=2$, the two beta coefficients in \eqref{eq:Q93:6.1}
are both $d$; for $a\ge3$, use $\beta\ge\kappa d$.
The sole $G_2$ exception is $a=2$, when $\|G_2\|\le6/5$;
all $G_a$ with $a\ge3$ are contractions. These observations and
the individual means in Table~\ref{tab:oc:q93-individual-means}
give Table~\ref{tab:oc:q93-unit-cost}. Its mean bounds are
$(7/4+1/5-2)d=-d/20$,
$(5/6+1/5-2\kappa)d=-37d/240$,
$(3/5+1/30-2)d=-41d/30$, and
$(3/5+1/30-2\kappa)d=-133d/240$, respectively.
The first two costs are
$-2/5+17(63/40)^2=66833/1600$ and
$-37/30+17(29/20)^2=41411/1200$.
In the third row the loss is
$128/5+164/15-17(27/20)^2=6661/1200>5$.
\begin{table}[htbp]
\centering\small
\setlength{\tabcolsep}{5pt}
\caption{Unit connectors. The mean, norm and cost columns are upper bounds. The final column is the lower loss at the minimum permitted pair count; each additional pair adds $128d/5$.}\label{tab:oc:q93-unit-cost}
\begin{tabular}{lccccc}
\toprule
Range & $k_{\min}$ & $(c_1+\E N)/d$ & $\|V\|/\sqrt d$ & Cost$/d$ & Loss$/d$\\
\midrule
$b_1=3,\ a=2$ & $3$ & $-1/20$ & $63/40$ & $66833/1600$ & $35487/1600$\\
$b_1=4,\ a\in\{2,3\}$ & $2$ & $-37/240$ & $29/20$ & $41411/1200$ & $20509/1200$\\
$b_1\ge5,\ a=2$ & $1$ & $-41/30$ & $27/20$ & $-164/15+17(27/20)^2$ & $6661/1200$\\
$b_1\ge5,\ a\ge3$ & $1$ & $-133/240$ & $23/20$ & $-133/30+17(23/20)^2$ & $9061/1200$\\
\bottomrule
\end{tabular}
\end{table}
For $a\ge3$ the last-row loss is
\begin{equation}\tag{Q93.6.2}\label{eq:Q93:6.2}
[128/5+133/30-17(23/20)^2]d=9061d/1200>d.
\end{equation}
Every further pair adds the same \(128d/5\) loss in all these cases.
Together with the nonunit calculation, this proves all three rows
of \eqref{eq:Q93:0.3}. Positive comparison restores \(E_t,E_r\).
Finally \(e^{-d}\le1/(1+d)\le1-3d/4\) because \(d\le1/3\).
Replacing \(K\) by \(-K\) reverses the receiving unit arm, inward arm,
all pairs, connector and decorations together while preserving the
measure and every \(\delta_i\).
